% Book template for D&D 5e adventures and sourcebooks
% Uses dndbook document class with structured content support

% Base LaTeX template for D&D 5e documents using DND-5e-LaTeX-Template
% This template provides the foundation for all D&D-style documents with structured content

\documentclass[letterpaper, 11pt, twocolumn]{dndbook}

% DND-5e-LaTeX-Template provides most styling automatically
% Only include additional packages if needed

% Additional packages for enhanced functionality
\usepackage{hyperref}
\usepackage{bookmark}
\usepackage{makeidx}

% Additional packages for book layout
\usepackage{multicol}
\usepackage{wrapfig}
\usepackage{tikz}
\usepackage{tcolorbox}
\usepackage{titlesec}

% DND book-specific environments
\tcbuselibrary{breakable}

\title{Dragon Delves: A Copper for a Song}
\date{\today}

% Disable automatic section numbering while preserving TOC entries
\setcounter{secnumdepth}{-1}

% Disable text justification
\raggedright
\raggedbottom

% Hyperlink configuration
\hypersetup{
    colorlinks=true,
    linkcolor=black,
    urlcolor=black,
    citecolor=black,
    filecolor=black,
    bookmarks=true,
    bookmarksopen=true,
    bookmarksnumbered=true,
    pdfsubject={D&D 5e Content},
    pdfkeywords={D&D, 5e, RPG, tabletop, }
}

% Enhanced chapter and section styling for books

% Configure section numbering depth for books
\setcounter{secnumdepth}{-1}

\begin{document}

\frontmatter

% Title page
\maketitle
\thispagestyle{empty}
\clearpage

% Table of Contents in frontmatter
\tableofcontents
\clearpage

\mainmatter

\chapter{Introduction}\label{ch:introduction-1-1}

% Image placeholder
\section{Dragons in Dungeons}

Nothing is more quintessentially Dungeons \& Dragons than gathering a group of heroes to face a mighty dragon in a mysterious dungeon. These classic elements link players across time and space, from the first adventurers who forayed into the gloom more than 50 years ago to players beginning their first adventure now.

\textit{Dragon Delves} is a celebration of this shared tradition. Each of the ten adventures in this book features a dragon in a lair. Designed for characters of levels 1 through 12, the adventures offer a range of challenges, from defeating evil dragons to working with good dragons, as well as navigating more complicated moral dilemmas. The dungeons where the dragons make their lairs range from cavern complexes to ruined cities and more.

Where there are dragons in dungeons, there's always adventure!
\section{Using This Book}

The adventures in an anthology such as \textit{Dragon Delves} are versatile resources that can fill a variety of needs at your gaming table: as one-shot adventures, as "filler episodes" in an ongoing campaign, or as a series that forms a campaign, bringing characters from level 1 to a climactic conclusion at level 12.

These adventures are designed to require minimal preparation (see area Running the Adventures later in this introduction).

\subsection{One-Shot Adventure}

Maybe you're normally a player in an ongoing campaign, but you want to try your hand at being the Dungeon Master. Maybe you're trying to pull a new group together, and you want to see how all the players get along before you start a whole campaign. Maybe your regular DM is sick or has a scheduling conflict, but everyone else is ready to play. Or maybe your group needs a break or a diversion from the ongoing campaign. In any of these cases, you can pick an adventure from this book and run it.

\subsection{In an Ongoing Campaign}

Maybe your players' characters know what they have to do to foil the evil boss's plans, but they're not quite high enough level to accomplish their goal. Or they've just hit a significant milestone in the campaign, and you want to give them a change of pace before they get immersed in the next act of the unfolding drama. Or they're traveling halfway around the world, and you're looking for something to keep the journey interesting---a little adventure along the way. Or you're not prepared for the next adventure, and you need a diversion! In any of these cases, you can pick an adventure from this book to serve as a "filler" episode or side quest in the ongoing campaign.

% Image placeholder

When the characters are the appropriate level for the adventure you wish to run, simply narrate how the characters end up in the adventure's locale. Each adventure's beginning includes a simple setting description, such as "a village with a nearby forest." Beyond these basic details, feel free to alter the adventure's location to match your campaign.

But if you do have time to prepare, you can incorporate the adventure you're running into the themes and threats of your campaign. You can alter the names of people and places in the adventure to fit better into your world, link the adventure villains to organizations and events that are important in your campaign, and plant clues to the themes of your campaign in the dragon's treasure hoard.
\section{Creating a Campaign}

These adventures can be strung together as a complete dragon-themed campaign. Simply guide your players through the adventures in the order presented in this book and award story-based level advancement at each adventure's end, so their characters are the appropriate level for the next challenge they'll face. (No level advancement is necessary before the final adventure.) You can also use one of the following narrative frameworks to connect the adventures to one another.

\subsection{Option 1: Mysterious Patron}

Have a friendly patron whose true identity is somewhat mysterious befriend the adventurers early in their careers and send them on each adventure in sequence. The patron might be a charismatic young monk, a kindly teacher, a diminutive scamp, or a scatter-brained elder, and might or might not be accompanied by seven golden canaries or dazzling yellow butterflies, which are actually ancient gold dragons in disguise---the companions of the dragon god Bahamut, which is the patron's true identity.

Perhaps foreseeing that they are destined to become mighty heroes, Bahamut takes the adventurers under his metaphorical wing and sends them on adventures to deal with draconic concerns, great and small, across the world. The patron might feign surprise when an adventure that doesn't at first seem to involve a dragon (such as "The Will of Orcus (p. 1)," chapter 3) turns out to revolve around one, but of course that's exactly the reason Bahamut sent the characters to investigate.

You can use each adventure as written, with the characters' patron steering them in the general direction of the adventure hook without spoiling any surprises in the adventure ahead. Alternatively, the patron could let the characters in on the secret early on, telling them that dragons will be a particular focus of their adventures and altering some of the adventure hooks along the way.

\subsection{Option 2: Brazen Egg}

Begin the campaign with the introductory encounter of chapter 10, "Dragons of the Sandstone City (p. 1)," ending with Eldemere entrusting the Brazen Egg to the characters. Eldemere's instructions are different, though: she explains that the egg must be "steeped" in the hoards of ten dragons, one of each chromatic and metallic kind, to activate its power. She also suggests that the egg will lead the characters toward suitable hoards. In fact, any character who sleeps within 15 feet of the egg has vivid dreams leading them toward the start of the next adventure.

To draw power from a dragon's hoard, the egg must spend 8 hours in the midst of the treasure, covered with coins and surrounded by other items. However, as soon as the empowered egg is brought within Anthradusk's hoard chamber (area H16 (p. 1) in chapter 10), it reaches its full potential and begins to crack open, as described in that adventure.

\subsection{Option 3: Hoard Magic Items}

If you have Fizban's Treasury of Dragons, p. 4, the hoard of magic items described in that book can provide an incentive for characters to seek out dragons of increasing power and different kinds. Add the following items to the rewards acquired in the first three adventures: a \textit{Scaled Ornament} in "Death at Sunset (p. 1)," a \textit{Dragon Vessel} in "Baker's Doesn't (p. 1)," and a \textit{Dragon's Wrath Weapon} and a \textit{Dragon-Touched Focus} in "The Will of Orcus (p. 1)." All four items are Slumbering when the characters first acquire them, but in later adventures, the items can be awakened by steeping them in dragons' hoards.
\section{Running the Adventures}

To run each of these adventures, you need the fifth edition core rulebooks: the Player's Handbook, Dungeon Master's Guide, and Monster Manual. Spells and equipment mentioned in the adventures are described in the Player's Handbook. Magic items are detailed in the Dungeon Master's Guide, and monsters appear in the Monster Manual.

The table of contents summarizes the adventures in this anthology. Each adventure is designed for four to six characters of a particular level, but you can adjust for larger or smaller groups by changing the number of foes in an encounter and giving a small group additional resources, as explained in the Dungeon Master's Guide. Each of the adventures can take place in any world you choose, as long as the location has the basic elements listed at the beginning of each adventure.

\subsection{Preparation}

Every adventure in this book begins with a page that helps you prepare and run the adventure. It includes the following elements:

\begin{itemize}
\item \textbf{Key Plot Points.} A summary of the key elements in the adventure's story
\item \textbf{Preparation.} Steps to get yourself ready to run the adventure, including a list of all the stat blocks from the Monster Manual you might need
\item \textbf{Key NPCs.} A table summarizing key nonplayer characters in the adventure---their names, roles, stat blocks, and locations---to help you keep track of them during play
\end{itemize}

\subsection{Player Assistance}

As described in the Dungeon Master's Guide, you can streamline your role as Dungeon Master by delegating certain tasks to the other players, especially if this is your first time behind the DM's screen. Before play begins, work with the other players to divide any or all of the following tasks among them:

\begin{itemize}
\item \textbf{Initiative Tracker.} In combat, one player tracks Initiative for the characters and monsters.
\item \textbf{Monster Wranglers.} Whenever the characters enter combat with Hostile monsters, one or two players run these monsters using the stat blocks in the Monster Manual. (Two players can share this role.)
\item \textbf{Rules Consultant.} One player references rules in the Player's Handbook or other books as needed.
\end{itemize}

\subsection{Adventures for One Character}

While the Dungeon Master's Guide explains how to adjust adventures for different party sizes, three of the adventures (indicated in the table of contents) are particularly well-suited for play with a single adventurer. Each of these adventures includes instructions for ensuring the character faces an appropriate challenge.

Central to those instructions is the \textit{Blessing of the Lone Champion}, a supernatural gift that bolsters the solo character's abilities.

\subsection{Blessing of the Lone Champion}

The blessing lasts for the duration of the adventure.

\begin{DndSidebar}{A Draconic History}
In addition to ten adventures, this book includes special sections that feature art of each dragon variety, drawing from almost the whole 50-year history of Dungeons \& Dragons, from the 1977 Monster Manual to the one published in 2025. The appearance of all ten dragon types has evolved significantly over the past fifty years, but close examination of these images reveals points of continuity as well, from the blue dragon's nose horn to the silver dragon's crest.
Artist Todd Lockwood spearheaded a thorough revision of all ten dragon types for the release of D\&D's third edition in 2000, basing his designs on the anatomy of real-world animals and certain points of continuity with earlier designs. With a few variations, those designs set the standard for D\&D's dragons for over 20 years.
In preparation for the release of the Monster Manual (2025), art director Josh Herman led a comprehensive refresh of all the dragon looks as D\&D began its second half-century. With the goal of ensuring each dragon's distinct identity---its personality, preferred habitat, and combat style---shines through in its look, this redesign showcases visually what's unique about each type of dragon.
In the same way, this book highlights how an adventure featuring one kind of dragon can feel distinct from all the others.
\end{DndSidebar}
\section{Adventures}

\begin{itemize}
\item Death at Sunset (Level 1) (p. 1)
\item Baker's Doesn't (Level 3) (p. 1)
\item The Will of Orcus (Level 4) (p. 1)
\item For Whom the Void Calls (Level 5) (p. 1)
\item The Dragon of Najkir (Level 7) (p. 1)
\item The Forbidden Vale (Level 9) (p. 1)
\item Before the Storm (Level 10) (p. 1)
\item Shivering Death (Level 11) (p. 1)
\item A Copper for a Song (Level 12) (p. 1)
\item Dragons of the Sandstone City (Level 12) (p. 1)
\end{itemize}

\chapter{A Copper for a Song}\label{ch:a-copper-for-a-song-2-2}

\begin{itemize}
\item \textit{ \textbf{An adventure for Level 12 characters.}}
\item \textit{This adventure is designed to fill one or two sessions of play.}
\item \textit{It can take place in any coastal town with nearby farmland and hills.}
\end{itemize}
This adventure takes place in Godsbreath, a D\&D setting introduced in Journeys Through the Radiant Citadel. You don't need that book to run this adventure, however; you can place it in any similar setting.
% Image placeholder
\section{Key Plot Points}

The following information is key to the adventure's story:

\begin{itemize}
\item \textbf{Deciphering Clues.} The adventurers undertake a quest to find the lost verses of a magical song that can rejuvenate the soil of a dying land. The characters must decipher clues hidden in another song's lyrics with the help of chatty locals.
\item \textbf{Braving the Rattle.} Having learned where the lost verses are, the characters travel to a copper dragon's lair in the hills of a region called the Rattle.
\item \textbf{Navigating Nakari's Lair.} The characters must convince Nakari, the copper dragon, to hand over the verses. Failing that, they can try to take the verses by force or learn them from a source hidden within the dragon's lair.
\end{itemize}
\section{Preparation}

Before running the adventure, prepare as follows:

\begin{itemize}
\item \textbf{Step 1.} If you're running the adventure for a single adventurer, read the "Running for One Character" section.
\item \textbf{Step 2.} Read "area Adventure Background" for more context and familiarize yourself with the Key NPCs (p. 1) table, which describes important individuals the characters will interact with throughout the adventure.
\item \textbf{Step 3.} Make a copy of the Copper for a Song handout for each player. (You will share it with them shortly after the adventure starts.)
\item \textbf{Step 4.} Bookmark the following stat blocks in the Monster Manual or on D\&D Beyond:
\end{itemize}

\begin{itemize}
\item \textbf{Axe Beak, Giant}
\item \textbf{Commoner}
\item \textbf{Copper Dragon, Adult}
\item \textbf{Djinni}
\item \textbf{Druid}
\item \textbf{Galeb Duhr}
\item \textbf{Green Dragon Wyrmling}
\item \textbf{Hill Giant}
\item \textbf{Kobold, Winged}
\item \textbf{Performer}
\item \textbf{Performer Maestro}
\item \textbf{Priest Acolyte}
\item \textbf{Purple Worm}
\item \textbf{Warrior Veteran}
\end{itemize}

% Table: Key NPCs
\begin{DndTable}[header={Key NPCs}]{llll}
Name & Role & Stat Block & Location \\
area Aunt Dellie & Human boat pilot & \textbf{Commoner} (Medium, Chaotic Good) & "area Listening Post" \\
Gleem Happenbolt & Dead gnome miner & --- & area N14 \\
area Hurch Henley & Human street magician & \textbf{Performer} (Small) & "area Listening Post" \\
area Kelra Ironweaver & Dwarf blacksmith & \textbf{Commoner} (Medium, Chaotic Neutral) & "area Listening Post" \\
Lashlo & Impatient guest of Nakari & \textbf{Djinni} & area N10 \\
Nakari & Keeper of the lost verses & \textbf{Adult Copper Dragon} & area N17 \\
Retta and Rynn & Wyrmlings adopted by Nakari & \textbf{Green Dragon Wyrmling} (Neutral) & area N17 \\
area Taelyn & Elf soldier & \textbf{Warrior Veteran} (Medium, Neutral Good) & "area Listening Post" \\
Tungsten Ward & Human quest giver & \textbf{Priest Acolyte} (Medium, Lawful Good) & "area Listening Post" \\
Yemi Moonsybil & Potential human ally under the effects of a \textit{Sequester} spell & \textbf{Performer Maestro} (Medium, Chaotic Good) & area N12 \\
area Zakia & Human herbalist & \textbf{Druid} (Medium, Chaotic Good) & "area Listening Post" \\
\end{DndTable}
\section{Running for One Character}

You can run this adventure for a single character of level 12. A Bard is a good choice, given the role that music and storytelling play in the adventure. Ideally, the character should have proficiency in at least two of the following skills:

\begin{itemize}
\item Deception
\item Intimidation
\item Perception
\item Performance
\item Persuasion
\item Stealth
\end{itemize}

The solo character gains the Blessing of the Lone Champion (see the introduction) to compensate for the lack of companions.
\section{Adventure Background}

The region of Godsbreath is known for its crimson earth, but the soil has grown depleted in recent years. Even as locals develop new ways to farm, the harvest produces less food each year, and the region grows closer to crisis.

Many believe that the key to restoring the soil lies in the verses of the \textit{Awakening Song}, a collection that forms a record of the region's history. Some of the verses have been lost for generations. Recently, an acolyte named Proclaimer Tungsten Ward learned from an old man about another acolyte who wrote down the verses of the \textit{Awakening Song} to preserve them. But the old man couldn't remember the acolyte's name and doesn't know what became of the written verses.

\subsection{The Verses}

Unknown to the people of Godsbreath, the lost verses now belong to an adult copper dragon named Nakari. She has lived in the rugged hills of the Rattle, an unsettled region north of Godsbreath, since before the first human settlers arrived to till the crimson soil.

Nakari liked to visit the settlers in human guise and learn their stories, songs, and ways. When she learned that an acolyte had dutifully written down the verses of the \textit{Awakening Song}, Nakari convinced the acolyte to sell her the record for a single copper coin and the promise that she would never let them fall into the wrong hands.
\section{Beginning the Adventure}

When you're ready to begin the adventure, read or paraphrase the following:

\begin{DndReadAloud}
A long voyage has brought you to Nightwater Cove. The town of Promise rises from the shore, its brightly painted buildings a welcome sight. The vessel docks, allowing you and the other passengers to disembark.The bustling harbor town rests on the southern coast of Godsbreath, a region known for its agriculture. Sadly, its farmlands have become depleted in recent years. The struggle to survive grows ever more dire, and this year's harvest is expected to be the poorest in recent memory. Locals believe they could rejuvenate the soil with the Awakening Song, whose ancient verses recount the history of the land and its people, but key verses have been lost.Several weeks ago, you received an invitation from Proclaimer Tungsten Ward, an acolyte who works at a temple in Promise called the Listening Post. Proclaimer Ward hopes you can find the lost verses of the Awakening Song, giving the people of Godsbreath the means to rejuvenate their dying soil.
\end{DndReadAloud}

When the characters head to the temple, proceed to "area Listening Post."
\section{Listening Post}

When the characters arrive at the area Listening Post, describe it as follows:

\begin{DndReadAloud}
The Listening Post is both a gathering place and a temple. Inside its colorful walls, small crowds are gathered around storytelling booths, a community bulletin board, and a trading hub. A human near the far wall waves at you and begins weaving through the crowd with scrolls tucked under one arm. You recognize Proclaimer Tungsten Ward. Slight of stature and neatly dressed in scholar's garb of purple and gold, Proclaimer Ward greets you in a soft voice and says, "The gods have delivered you safely. For that and many other things, I am grateful." They then hand each of you a scroll.
\end{DndReadAloud}

Proclaimer Tungsten Ward (Medium, Lawful Good \textbf{Priest Acolyte}) gives each character a nonmagical scroll with lyrics written on it in Common.

Give each player a copy of the Copper for a Song handout at this time and invite one or more of them to read the song lyrics aloud.

% Image placeholder

\subsection{Proclaimer Ward's Quest}

When the characters are ready to learn more from Proclaimer Ward, read or paraphrase the following:

\begin{DndReadAloud}
"Stories and songs are the lifeblood of Godsbreath," says the young acolyte, "and the Awakening Song is woven from the stories of Godsbreath's first settlers and their descendants. "Since I contacted you about my quest to find the lost verses of the Awakening Song, I learned that another acolyte had filled a book with the song's lost verses. This happened so long ago, the acolyte is surely dead and mostly forgotten. "The scroll I gave you holds 'Copper for a Song,' a song they wrote describing an encounter with a dragon named Nakari. Within these lyrics are clues to finding the lost verses of the Awakening Song, I'm sure of it. "I've gathered people here to help decipher the clues in 'Copper for a Song.' We need your help. And when we discover where the verses are, I'll need you to retrieve them."
\end{DndReadAloud}

\subsubsection{Reward}

If the characters demand a reward, Proclaimer Ward offers to give them a \textit{Figurine of Wondrous Power (onyx dog)} once they recover the lost verses of the \textit{Awakening Song}. The figurine occasionally serves as a temple guardian, but Proclaimer Ward is willing to part with it.

\subsection{Listening Post Conversations}

Proclaimer Ward steers the characters to a small group of other people in the area Listening Post to workshop the clues buried in the "Copper for a Song" lyrics. Proclaimer Ward introduces them to the five other group members and provides basic details about each person, described below.

Each of these five individuals is initially Indifferent toward the characters. A character who takes an Influence action and strikes up a conversation with one of them can make a DC 14 Charisma (Persuasion) check to try to shift the person's attitude from Indifferent to Friendly. Alternatively, a character can try to improve a person's attitude with a different ability check. For example, a character can attempt a Dexterity (Sleight of Hand) check to impress area Hurch with a coin trick or an Intelligence (Nature) check to impress area Zakia with knowledge about rare herbs or monsters.

A person who is Friendly toward a character helps decipher clues in the song lyrics and imparts the information below when shown the scroll.

\subsubsection{Aunt Dellie}

\textbf{Aunt Dellie} (Medium, Chaotic Good \textbf{Commoner}) is a human skimmer boat pilot with a skeptical mind, a soft heart, a loud laugh, and a talent for learning a bit about everyone's business. She wears a simple yellow dress and keeps her hair short.

When Friendly, Aunt Dellie shares the following: "Past the farms of the Ribbon is a wild, unsettled land called the Rattle. We sometimes call its bald hills 'hills of gold' because of their color. The phrase 'beyond the settled land' must refer to the Rattle---a land of monsters."

\subsubsection{Hurch Henley}

\textbf{Hurch Henley} (Small \textbf{Performer}) is a human street magician whose attempts at card tricks and legerdemain are less interesting than his stories from when he worked on his father's farm. He wears a floppy straw hat and a loose vest.

When Friendly, Hurch shares the following: "Father sold a few pigs to a copper-haired witch named Nakari. She said she lived four miles inside the Rattle, but Father thought she was joking. 'The Rattle is too evil a place for people to live,' he used to tell me."

\subsubsection{Kelra Ironweaver}

\textbf{Kelra Ironweaver} (Medium, Chaotic Neutral \textbf{Commoner}) is a local dwarf blacksmith who has outlived five husbands and knows a lot about history because she was there when it happened. Her hair cascades in thick ringlets over her broad shoulders.

When Friendly, Kelra shares the following: "My first husband, Irvin, was a copper miner in the Rattle. He claimed a dragon lived under a hill not far from the mine. Poor Irvin---flattened by a hill giant!"

\subsubsection{Taelyn}

\textbf{Taelyn} (Medium, Neutral Good \textbf{Warrior Veteran}) is an elf soldier with a Tiny blue finch perched on the right shoulder of her green cloak. She barters for supplies at the area Listening Post's trading hub.

When Friendly, Taelyn shares the following: "Last spring, a bulette terrorized farms at the north end of the Ribbon. One day, out of the blue, a copper dragon swooped down and dissolved the bulette with its acid breath, saving a farmer's life."

\subsubsection{Zakia}

\textbf{Zakia} (Medium, Chaotic Good \textbf{Druid}) is a local human herbalist who uses one of the storytelling booths to spook young children with tales about dragons and child-eating monsters. Her blue dress is stained with dirt from her gardens, and a bag she wears slung over her shoulder is overflowing with aromatic plants.

When Friendly, Zakia shares the following: "If a copper dragon is what you seek, that's mostly good news. Copper dragons are even-tempered but also miserly. They're as fond of music, jokes, and riddles as they are of treasure. Their burrows are often sinuous, winding, and riddled with secret doors."

% Image placeholder

\subsection{Putting It Together}

After conversing with the townsfolk in the area Listening Post, the characters should have enough information to find the dragon's lair, which is located in the Rattle, 4 miles past the Ribbon, under a hill near an old copper mine.
\section{The Rattle}



Godsbreath map shows the location of Nakari's lair, roughly 12 miles northeast of Promise in the Rattle as the dragon flies, or 20 miles along the winding northern road. The area is shaped by the presence of a copper dragon's lair, as described in the Monster Manual.

\subsection{Danger in the Hills}

% Image placeholder

To get to the Rattle, the characters must cross the stretch of farmland called the Ribbon.

Read or paraphrase the following to describe the journey into the Rattle:

\begin{DndReadAloud}
The journey from Promise leads north through the Ribbon, where farmhouses dot the dark-red ground between stands of trees. Eventually, the road dwindles to a track. The farms grow fewer and farther between as you approach the Rattle.After about four miles, the trees give way to rocky hills of a golden-brown hue. The entrance of an old copper mine is blocked with boulders and serves as a landmark. Sprawled nearby is a huge, partly dissolved monster with its guts ripped out. The dead monster lies on its armored back with its four clawed feet in the air. From deeper in the hills come the sounds of loud belching and guttural laughter.
\end{DndReadAloud}

The copper mine contains nothing of interest. A character can take the Study action and make a DC 12 Intelligence (Nature) check to identify the dead monster, recognizing it as a bulette on a success.

As a Search action, a character can inspect the dead bulette and make a DC 12 Wisdom (Medicine) check. On a success, the character determines that the creature died from intense acid and then was partially eaten by an enormous predator.

After finding the bulette, the characters might follow the sounds of belching and laughter (see "area Hill Giant Camp") or might reconnoiter the area (see "area Tunnels to the Dragon's Lair"). Use the "Surface" half of Nakari's Lair (p. 1) map for either approach.

\subsubsection{Hill Giant Camp}

If the characters follow the sounds coming from the hills, read or paraphrase the following:

\begin{DndReadAloud}
Six hill giants relax in a clearing between rocky outcroppings. Crystalline deposits in the rock glitter brightly, but the giants seem oblivious to the natural beauty of their surroundings. They sit in a rough circle, belching and laughing loudly. Two equally enormous flightless birds with axe-like beaks shift restlessly nearby.Several of the outcroppings have ten-foot-wide tunnels bored into them. The giants pay no attention to these tunnels.
\end{DndReadAloud}

Use the "Surface" half of Nakari's Lair (p. 1) map for this encounter. The characters approach from the east.

The glittering crystalline formations in the rock drew Nakari to this area. These crystals are pretty but not valuable.

\paragraph{Hill Giants}

The six \textbf{Hill Giants} are Hostile and unaware that they're camping above the subterranean lair of a copper dragon. Their two \textbf{Giant Axe Beaks} serve as beasts of burden and hunting companions, and they fight at the giants' command.

As an Influence action, a character can attempt one of the following:

\begin{itemize}
\item \textbf{Bribery.} A character can try to befriend a giant by giving it at least 10 pounds of food, 5 gallons of water or ale, or 50 GP worth of treasure. The character makes a DC 13 Charisma (Persuasion) check. On a success, the attitude of the giant shifts to Indifferent. The other giants demand similar payment; if the character refuses, the Hostile giants begin combat. On a failure, the Hostile giants begin combat.
\item \textbf{Fear.} A character who knows Giant can try to frighten the giants by telling them about the dragon who lives nearby. The character makes a DC 16 Charisma (Intimidation) check. On a success, one pair of giants leaves and never returns. On a failure, Hostile giants begin combat.
\item \textbf{Performance.} A character can try to amuse the giants with music, dancing, juggling, capering, loud belching, or some other performance. The character makes a DC 13 Charisma (Performance) check. If the check succeeds, the giants' attitudes toward that character shift to Friendly, and their attitudes toward the other characters shift to Indifferent. On a failure, Hostile giants begin combat.
\end{itemize}

\paragraph{Treasure}

The six giants share three enormous bags to store their loot:

\begin{itemize}
\item \textbf{Sack 1}
\item \begin{itemize}
\item 35 GP, 270 SP, and 1,800 CP
\item Five wheels of stinky cheese
\item Dire wolf's skull
\end{itemize}
\item \textbf{Sack 2}
\item \begin{itemize}
\item Rusty Dagger (which the giants use as a toothpick)
\item 180 SP and 2,200 CP
\item Unopened cask of honey mead (worth 50 GP)
\item Beehive with live bees
\end{itemize}
\item \textbf{Sack 3}
\item \begin{itemize}
\item 150 SP and 2,500 CP
\item Haunch of smoked meat
\item Smelly, giant-size sandal
\item \textit{+2 Shield}
\end{itemize}
\end{itemize}

\subsubsection{Tunnels to the Dragon's Lair}

Characters who reconnoiter the area notice four tunnel entrances (marked on the "Surface" half of Map: Nakari's Lair (p. 1)), all of which lead down to the dragon's lair. The tunnels are too small for the giants to enter comfortably, so they don't even try.

The entrances to tunnels area N1 and area N2 lie within the giants' line of sight; reaching either tunnel without being noticed requires a successful DC 13 Dexterity (Stealth) check. Rocky outcroppings hide the entrances to tunnels area N3 and area N4 from the giants' view.
\section{Nakari's Lair}

Nakari is in her lair when the characters arrive, raising a pair of green dragon wyrmlings she recently adopted; see her den (area area N17) for details. Nakari is initially Indifferent toward intruders.

\subsection{Lair Features}

Nakari's lair is a network of tunnels and caves hewn from solid rock. Other features of the lair are summarized below.

\subsubsection{Ambient Music}

While Nakari is in her lair, soft music permeates the area. The music, which blends woodwinds and bells, is faint and has no discernible source.

\subsubsection{Breakaway Roofs}

Where denoted in area descriptions, certain parts of the lair have 5-foot-thick breakaway roofs that any creature with a Burrow Speed can break through to create an opening to the surface. As a Search action, a character who examines a breakaway roof can make a DC 15 Intelligence (Investigation) check. On a success, the character can tell that the roof has been structurally weakened but is in no danger of collapsing on its own. A character with Tremorsense has Advantage on this check.

\subsubsection{Ceilings}

Ceilings are 10 feet high throughout the lair.

\subsubsection{Laughing Gas}

Areas tinted purple on Nakari's Lair (p. 1) map are filled with magical laughing gas that is visually imperceptible but smells like candy. Any creature other than Nakari that enters the gas-filled area for the first time on a turn or starts its turn there must succeed on a DC 15 Wisdom saving throw or have the \textit{Incapacitated} condition until the start of its next turn, as it laughs uncontrollably.

A strong wind (such as that created by \textit{Gust of Wind}) disperses the gas for 1 minute, while \textit{Dispel Magic} disperses the gas for 24 hours. The gas dissipates if Nakari dies.

\subsubsection{Lighting}

The lair is in Darkness. Nakari and other denizens rely on Darkvision to see. Descriptions assume the characters have a way of seeing in the dark.

\subsubsection{Secret Doors}

All secret doors in Nakari's lair are magical and marked on the map. Each secret door blends seamlessly with the surrounding wall (DC 20 to notice). Nakari can open any secret door as a Bonus Action from anywhere in her lair, causing the stone to reshape into an opening wide enough for her to pass through. Other creatures must break through the stone secret door. Whenever a creature with an Intelligence of 3 or higher comes within 10 feet of a secret door, Nakari becomes aware of the creature's presence and location but takes no immediate action.

\subsection{Lair Locations}



The following locations are keyed to the "Underground" half of Map: Nakari's Lair.

\subsubsection{N1: East Entrance}

\begin{DndReadAloud}
Ethereal music travels up this dark, sloping tunnel.
\end{DndReadAloud}

This tunnel slopes down to area area N6. A branching tunnel leads to area area N5.

\subsubsection{N2: South Entrance}

This area is identical to the east entrance (area area N1), except the tunnel slopes down to area area N8.

\subsubsection{N3: North Entrance}

This area is identical to the east entrance (area area N1), except the tunnel slopes down to area area N11.

\subsubsection{N4: West Entrance}

This area is identical to the east entrance (area area N1), except the tunnel slopes down to area area N13.

\subsubsection{N5: Eastern Boulder Cave}

\begin{DndReadAloud}
Several boulders are scattered about this small cave, and three of them raise their heads to look at you as you enter. Crooked grins full of rocky teeth appear in their faces as the creatures rumble what might be a greeting.
\end{DndReadAloud}

Three gregarious \textbf{Galeb Duhr} dwell in this cave. They are pleased to have company, and ask in Primordial if the characters want to join them in a game of rock tossing. If the characters don't understand their words, the galeb duhr try to get the idea across through pantomime. The goal of the game is to hurl enormous, 200-pound rocks to land as close as possible to a target on the ground. If the characters refuse to play, the galeb duhr are disappointed and return to their rest.

\paragraph{Development}

The \textbf{Hill Giants} outside the dragon's lair notice the sounds of rock tossing or combat in this cave. If the giants aren't aware of the characters' presence, one giant investigates, situating itself outside the tunnel to area area N1 for 10 minutes to find out if anything emerges.

\subsubsection{N6: Kobold Cave}

As the characters approach this cave, high-pitched chatter issues from within. Characters who know Draconic discern the following phrases in that language:

\begin{itemize}
\item "Nakari will be pleased."
\item "Too much clay! The wings will break!"
\item "The tail is too long."
\end{itemize}

The cave inhabitants stop talking if the characters have a light source with them or otherwise don't mask their approach.

When the characters enter the cave, read or paraphrase the following:

\begin{DndReadAloud}
Seven kobolds with leathery wings are making a statue out of clay in the middle of this cave, which has exits to the northwest, south, and east. The statue, about five feet tall and eight feet long, bears a passing resemblance to a dragon.
\end{DndReadAloud}

The seven \textbf{Winged Kobolds} want to impress Nakari with their artistic skills and thus win her favor. The kobolds' names are Drix, Grox, Kex, Mox, Nix, Prax, and Zox.

\paragraph{Roleplaying the Kobolds}

The kobolds are initially Indifferent toward the characters. One or more characters can shift the kobolds' attitudes to Friendly by helping them finish the statue (see "area Sculpting the Statue") or by giving them treasure or worthless trinkets that easily could be mistaken for treasure.

If the kobolds' attitude shifts to Friendly, the kobolds impart the following information in conversation with the characters:

\begin{itemize}
\item \textbf{Nakari.} The copper dragon, Nakari, loves art. The kobolds hope she will appreciate the clay statue they've crafted for her. The kobolds don't know where the dragon keeps her hoard; they've had no luck finding it.
\item \textbf{Secret Doors.} Nakari's lair contains secret doors that the kobolds have witnessed the dragon open magically. (The kobolds can lead characters to the secret doors concealing areas area N7, area N12, and area N15, but the kobolds can't open them. The kobolds are unaware of the secret door concealing area area N16, having never visited that part of the lair.)
\item \textbf{Wyrmlings.} Nakari rescued a pair of wyrmlings from the nearby forest and is raising them as her own. Both wyrmlings have tried to murder the kobolds more than once. (The kobolds don't know the wyrmlings' names.)
\end{itemize}

\paragraph{Sculpting the Statue}

The statue is meant to depict an adult copper dragon but needs a lot of work. A character who spends at least 10 minutes resculpting part of the statue makes a DC 15 Charisma (Performance) check. Each successful check improves one of four sections of the statue: the head, the wings, the body, or the tail. Once all four sections are improved, the kobolds are satisfied, praise the characters' artistic abilities, and become Friendly.

\paragraph{Treasure}

Each kobold has a Pouch that contains 1d12 CP.

\subsubsection{N7: Molting Cave}

This cave lies hidden behind a secret door (see "area Lair Features").

\begin{DndReadAloud}
This dark cave has no other exits. Scattered over the floor are molted copper dragon scales. Toward the back of the cave lie three small piles of treasure. Someone has carved a word in Draconic into the floor next to each pile.
\end{DndReadAloud}

Nakari keeps her shed scales here.

\paragraph{Breakaway Roof}

Nakari has carefully weakened the cave's roof so she can burrow through it quickly to reach the surface (see "Lair Features").

\paragraph{Treasure}

Nakari has gathered an assortment of baubles into three neat piles labeled in Draconic: "keep," "trade," and "dump." Many of these items were gifts from travelers or admirers; the remaining items the dragon acquired during her efforts to better understand and appreciate her neighbors.

The first time anyone other than Nakari removes an item from the "keep" or "trade" pile, the copper scales on the floor suddenly swirl through the air, slashing creatures with their sharp edges. Each creature in the cave makes a DC 15 Dexterity saving throw, taking 21 (6d6) Slashing damage on a failed save or half as much damage on a successful one. This magical effect can be triggered only once.

Each pile's contents are listed below:

\begin{itemize}
\item \textbf{Keep Pile}
\item Bagpipes (worth 30 GP)
\item \begin{itemize}
\item Ten Books of poetry (worth 25 GP each)
\item Twenty \textit{Holy Symbols} (amulets worth 5 GP each) representing various faiths
\item Framed painting (worth 7,500 GP) of a silver-haired, golden-eyed aasimar whose facial expression, attire, and hairstyle magically change daily at dawn
\item Eight music boxes (worth 250 GP each)
\item Twenty gold rings (worth 25 GP each)
\end{itemize}
\item \textbf{Trade Pile}
\item \begin{itemize}
\item Ten Books (historical texts worth 25 GP each)
\item Five flutes (worth 2 GP each)
\item Bejeweled, fur-lined cloak fit for royalty (worth 2,500 GP)
\item Ten gold rings (worth 25 GP each)
\item Fifty sheets of music (worth 1 GP each)
\end{itemize}
\item \textbf{Dump Pile}
\item \begin{itemize}
\item Five Books of poetry (worth 25 GP each)
\item Drum that's battered and worn (worth 1 GP)
\item Lute with two broken strings (worth 15 GP as is or 35 GP if repaired with a \textit{Mending} spell)
\item Three-dragon ante Gaming Set (worth 1 GP)
\item Ten gold bracelets (worth 25 GP each)
\item Ten silver rings (worth 5 GP each)
\item Rolled-up, torn portrait of a famous tiefling musician (worth 750 GP if repaired with a \textit{Mending} spell)
\end{itemize}
\end{itemize}

\subsubsection{N8: Southern Boulder Cave}

\begin{DndReadAloud}
Several boulders are scattered about this cave, and the ground vibrates slightly beneath you.
\end{DndReadAloud}

This cave appears to hold nine boulders, but three are \textbf{Galeb Duhr} that are singing while rolled up. Their song is quiet but deep and resonant, causing vibrations in the floor. If the galeb duhr notice the characters, they lift their heads and begin singing more loudly, shaking the ground and making bits of gravel fall from the ceiling as they try to entertain their unexpected audience. When the characters leave, the galeb duhr return to their quieter song.

\paragraph{Development}

The \textbf{Hill Giants} outside the dragon's lair notice the galeb duhr's loud singing. If the giants aren't aware of the characters' presence, one giant investigates, situating itself outside the tunnel to area N2 for 10 minutes to find out if anything emerges.

\subsubsection{N9: Laughing Cave}

\begin{DndReadAloud}
This empty cave smells like candy. It has exits to the west, north, and southeast.
\end{DndReadAloud}

\paragraph{Breakaway Roof}

Nakari has carefully weakened the roof of this cave so she can burrow through it quickly to reach the surface (see "area Lair Features").

\paragraph{Laughing Gas}

This area is filled with magical laughing gas (see "Lair Features").

\subsubsection{N10: Reception Cave}

\begin{DndReadAloud}
This oddly shaped cave smells like candy and has exits to the west and southeast. An alcove in the northeast corner is arrayed with colorful cushions, and a towering, purple-skinned woman in luxurious silk clothing lounges there, holding a golden chalice.
\end{DndReadAloud}

A peevish \textbf{Djinni} named Lashlo rests in the alcove. A guest that Nakari has left waiting far too long, Lashlo hoped to trade stories and songs with Nakari, but her patience is nearly at an end. She is Hostile to the characters, whom she considers rivals for Nakari's attention, and she fights them if they aren't effusively polite. If any character succumbs to the laughing gas for 2 consecutive rounds, Lashlo assumes they're laughing at her and fights them. The genie is unaffected by the laughing gas.

\paragraph{Laughing Gas}

This cave is filled with laughing gas (see "Lair Features").

\paragraph{Treasure}

The chalice belongs to Nakari, but the dragon won't notice its absence, and it's worth 750 GP. If the genie dies, her body dissolves into mist, leaving behind her silk clothing (worth 250 GP), two black pearl gemstones (worth 500 GP each), and a \textit{Wand of Lightning Bolts}.

\subsubsection{N11: Northern Boulder Cave}

\begin{DndReadAloud}
Several boulders lie scattered about this cave. As you enter, three of them raise their heads to look at you. They motion frantically for you to be quiet.
\end{DndReadAloud}

Three \textbf{Galeb Duhr} dwell in this cave. They are desperate to avoid the unwelcome attention of Lashlo, the djinni in area area N10, and try to remain as silent as possible.

\paragraph{Development}

The galeb duhr's caution is unwarranted. Noise in this cave doesn't attract the attention of the djinni or of the giants outside the dragon's lair.

\subsubsection{N12: Maestro's Cave}

This cave lies hidden behind a secret door (see "area Lair Features").

\begin{DndReadAloud}
This cave appears to be empty. At the back, two narrow tunnels wrap around a pillar of rock, leading to a smaller cave to the east. In the middle of that cave rests a slab of granite that measures three feet wide, eight feet long, and two feet high.
\end{DndReadAloud}

The western half of the cave is indeed empty (but see "area Breakaway Roof" below). An \textit{Invisible} creature lies on the stone slab in the smaller cave to the east (see "area Sequestered Maestro").

\paragraph{Breakaway Roof}

Nakari has carefully weakened the roof in the western half of this cave so she can burrow through it quickly to reach the surface (see "area Lair Features").

\paragraph{Sequestered Maestro}

Lying atop the stone slab is the human Yemi Moonsybil (Medium, Chaotic Good \textbf{Performer Maestro}). Yemi has the \textit{Invisible} condition and is in a state of suspended animation, both courtesy of a \textit{Sequester} spell. Nakari keeps dust and other detritus from accumulating on the slab and Yemi's resting form.

Yemi wears a glass key keepsake on a copper chain around her neck, and she gently clasps an \textit{Instrument of the Bards (Ollamh harp)}, which she acquired in her youth. Taking the harp, plucking one of its strings, or removing the glass key from around Yemi's neck ends the \textit{Sequester} spell early.

Yemi, who has been in stasis for almost a century, looks to be in her fifties. She knows all the lost verses of the Awakening Song.

Once she is no longer under the \textit{Sequester} spell's effect, Yemi is hungry for news. If the characters explain their quest, she sympathizes with the dire situation in Godsbreath and agrees to share the lost verses of the Awakening Song and return to Promise. However, Yemi won't leave Nakari's lair without first saying goodbye to the dragon. Yemi is familiar with the lair and can find her way around it; if Nakari is still alive, the dragon opens any secret door in Yemi's way.

\paragraph{Yemi's Story}

In her twenties and thirties, Yemi belonged to a subversive group of musicians called the Key. They used their music to inspire rebellions in lands plagued by tyrants and religious zealots---and made many powerful enemies in doing so. Decades later, after an assassin killed Yemi's spouse Mertyl Swooney in pursuit of the maestro, Yemi sought out Nakari. Sympathetic to Yemi's plight, Nakari revealed her true draconic nature and, with the maestro's consent, used a \textit{Spell Scroll} of \textit{Sequester} to make it impossible for enemies to find Yemi. After a century in hiding, Yemi figures her old enemies are long gone, and she's ready to return to the world.

% Image placeholder

\subsubsection{N13: Worm Cave}

\begin{DndReadAloud}
A pillar of rock braces the roof of this dry cave, and the ground rumbles under your feet.
\end{DndReadAloud}

A \textbf{Purple Worm} has been drawn to Nakari's lair by the dragon's frequent excavation work. The worm bursts up through the floor of the cave when the characters arrive and attacks them. It retreats when it becomes Bloodied.

\paragraph{Breakaway Roof}

Nakari has carefully weakened the roof of the cave so she can burrow through it quickly to reach the surface (see "area Lair Features").

\subsubsection{N14: Western Boulder Cave}

\begin{DndReadAloud}
Several boulders are scattered about this dark, sunken cave. Slumped in a niche in the southwest corner is a dead gnome clutching a well-made pick.
\end{DndReadAloud}

This cave appears to hold nine boulders, but three are belligerent \textbf{Galeb Duhr} relaxing while rolled up. If they notice the characters, they arise and begin combat.

\paragraph{Dead Gnome}

As a Search action, a character can examine the dead gnome and make a DC 12 Wisdom (Medicine) check. On a success, the character confirms that the gnome was battered to death roughly ten days ago. Characters who cast \textit{Speak with Dead} on the corpse can learn the following information:

\begin{itemize}
\item \textbf{Gnome's Name.} The gnome's name was Gleem Happenbolt.
\item \textbf{Home.} The gnome came from a mining settlement five miles north of the Rattle.
\item \textbf{War Pick.} The gnome's weapon is a family heirloom Gleem inherited from his uncle (see "Treasure" below).
\item \textbf{What Happened.} The gnome was searching the hills for gem deposits when he saw a copper dragon flying overhead. Believing the dragon to be friendly and thinking it might lead him to treasure, he followed it. The dragon led Gleem to this cave complex, but some of the boulders here "woke up" and tried to scare him away. When he wouldn't leave, they battered him to death.
\end{itemize}

\paragraph{Treasure}

A search of the dead gnome yields an \textit{Explorer's Pack} (with a full \textit{Waterskin}) and a \textit{+1 War Pick}. Characters also notice that the gnome's right boot has a pocket sewn into its side. The boot-pocket contains two topaz gemstones worth 500 GP each.

\subsubsection{N15: Tunnel to the Dragon's Den}

This tunnel is filled with laughing gas (see "area Lair Features").

\subsubsection{N16: Hidden Hoard}

This cave is hidden behind a secret door (see "area Lair Features"). If the characters attempt to break through or bypass this door, Nakari loudly invites them to join her in her den (area area N17)---and begins combat with them if they refuse.

\begin{DndReadAloud}
A sloping tunnel leads down to a cave containing a mound of copper coins sprinkled with silver, gold, and platinum coins. Stacked neatly against the wall behind the hoard are several paintings and other art objects.
\end{DndReadAloud}

\paragraph{Lost Verses}

Buried deep within the pile of coins is a ratty old notebook with a string tied around it. A character who spends 10 minutes searching through the coins finds the notebook.

Titled \textit{Verses of the Awakening Song}, the notebook was written by a young man who identifies himself as Proclaimer Cirroc Coldwater, Acolyte of the Covenant. The book is valuable only to the residents of Godsbreath, for it contains lost verses of the \textit{Awakening Song}. If the characters take it without Nakari's consent, the dragon does everything in her power to reclaim it.

\paragraph{Treasure}

In addition to the notebook containing the Awakening Song, Nakari's hoard contains the following treasure:

\begin{itemize}
\item 500 PP, 10,000 GP, 55,000 SP, and 450,000 CP
\item Four framed masterpiece paintings of brass, bronze, copper, and gold dragons (worth 2,500 GP each) by the renowned dragonborn painter Zaemon Vedrak
\item Pearl-studded, clamshell-shaped music box (worth 750 GP)
\item \textit{Sending Stones}
\item \textit{Talking Doll} that resembles a baby copper dragon
\end{itemize}

\subsubsection{N17: Nakari's Den}

\begin{DndReadAloud}
This cave has exits to the west and east, and its walls are patterned to look like dragon scales. In the middle of the cave sits a great copper dragon. Two green dragon wyrmlings crawl over and around her, occasionally exhaling puffs of gas from their nostrils.At the sight of you, the wyrmlings hiss and bare their fangs. The copper dragon smiles and says, "Pay them no mind. It's me you must answer to."
\end{DndReadAloud}

The \textbf{Adult Copper Dragon} is Nakari. The two \textbf{Green Dragon Wyrmlings} are Retta and Rynn. A little over a month ago, Nakari found the wyrmlings in the Rattle and thought it best to adopt them before they harmed themselves or someone else. The orphans are Neutral; they have bonded with their adoptive mother and don't know what happened to their biological parents.

\paragraph{Roleplaying Nakari}

Nakari is trying to teach the dragon wyrmlings to be kind toward others, so she adopts an Indifferent attitude toward the characters, even though she considers their intrusion unwelcome. Her attitude shifts to Hostile if one or more characters threaten the wyrmlings or disrespect her in her presence. Once her attitude shifts to Hostile, nothing the characters do from that point on can improve it.

Nakari doesn't like dealing with folk she doesn't know, and she distrusts anyone who professes to have noble intentions. She politely asks the characters to leave in peace and wishes them a safe journey home.

Nakari won't sell \textit{Verses of the Awakening Song} for any price, despite paying only a single Copper Piece for it.

\paragraph{Fighting Nakari}

If combat seems likely, Nakari sends the wyrmlings to their room (area area N18), where they remain until she calls for them. If Yemi accompanies the characters (see area area N12), Nakari won't do anything that might endanger her.

\paragraph{Breakaway Roof}

Nakari has carefully weakened the roof of the cave so she can burrow through it quickly to reach the surface (see "area Lair Features"). If battle erupts and the outcome seems likely to go against her, she breaks through the roof and orders the wyrmlings to flee by that route while she covers their escape.

\paragraph{Impressing Nakari}

As an Influence action, a character can recite the lyrics to "Copper for a Song (p. 1)," tell a story, or play an instrument, and then make a DC 15 Charisma (Performance) check. On a success, Nakari becomes Friendly and loans the characters her copy of Verses of the Awakening Song (see area area N16). The loan extends for up to 1 month, during which time the characters can have copies of the book made. If Yemi is with the party (see area area N12), Nakari entrusts the book to her; the maestro promises to return it within the agreed-on time frame.

\subsubsection{N18: Wyrmlings' Nursery}

\begin{DndReadAloud}
A short, downward-sloping tunnel leads to a dead-end cave. Its floor is littered with copper dragon scales, torn bales of hay, and broken toys.
\end{DndReadAloud}

If present, the two \textbf{Green Dragon Wyrmlings} Retta and Rynn are quick to lash out at intruders in this cave. The wyrmlings surrender or flee if either of them becomes Bloodied.
\section{Conclusion}

Characters who obtain the lost verses of the \textit{Awakening Song} can deliver the book safely to Proclaimer Tungsten Ward in Promise. Alternatively, they can accomplish their goal by escorting Yemi Moonsybil back to town. If either event comes to pass, word spreads quickly that the lost verses of the \textit{Awakening Song} have been recovered.

The townsfolk share the lost verses widely, ensuring that they're never forgotten again. The verses' lyrics tell a story about a previous struggle with harvests. From the details of how the locals solved that problem, Godsbreath's people coax the land back to peak fertility.

Characters who parted with Nakari on good terms can revisit the copper dragon from time to time. She might loan them some useful magic items (such as the \textit{Sending Stones} in her hoard) or hire them to complete a quest, such as finding out what happened to Rynn and Retta's parents.

% Image placeholder
\section{History of Copper Dragons}

RACH DARASTRIX

\begin{DndReadAloud}
Two long horns and prominent cheek frills are consistent elements of the copper dragon's look throughout D\&D's history. These features make it clear that Lord Gunthar is riding a copper dragon on Larry Elmore's cover for the 1985 adventure \textit{Dragons of Deceit}.
\end{DndReadAloud}

% Image placeholder: Dragons of Deceit (1985)

% Image placeholder: Monstrous Manual (1993)

% Image placeholder: Monster Manual (1977)

% Image placeholder: Draconomicon (2003)

% Image placeholder: Commander Legends: Battle for Baldur's Gate (2022)

% Image placeholder: Draconomicon (2003)

% Image placeholder: Monster Manual 2 (2009)

% Image placeholder: Monster Manual (2014)

% Image placeholder: Concept Art (2023)

% Image placeholder: Monster Manual (2025)

\chapter{Credits}\label{ch:credits-3-3}

\begin{itemize}
\item \textbf{Lead Designers:} Amanda Hamon, James Wyatt
\item \textbf{Designers:} Justice Ramin Arman, Makenzie De Armas, Anthony Joyce-Rivera, Renee Knipe, Ron Lundeen, Christopher Perkins, Patrick Renie, Erin Roberts, F. Wesley Schneider, Carl Sibley, Jason Tondro, Steve Townshend
\item \textbf{Rules Developers:} Makenzie De Armas (lead), Jeremy Crawford, Ron Lundeen
\item \textbf{Editors:} Judy Bauer (lead), Emma Candon, Sadie Lowry, Hannah Rose, Sue Weinlein, Hans Ziegler
\item \textbf{Art Directors:} Matt Cole (lead), Fury Galluzzi, Josh Herman, Kate Irwin, Emi Tanji
\item \textbf{Graphic Designers:} Matt Cole, Paolo Vacala
\item \textbf{Cover Illustrators:} Justine Jones, Greg Staples
\item \textbf{Interior Illustrators:} Tom Babbey, Daren Bader, Clyde Caldwell, Chippy, Simon Dominic, Jeff Easley, Luke Eidenschink, Larry Elmore, Jason A. Engle, Lars Grant-West, John Grello, Johan Grenier, Justine Jones, Andrew Kolb, Vance Kovacs, Andrey Kuzinskiy, Ed Kwong, Todd Lockwood, Antonio José Manzanedo, Dominik Mayer, Caio Monteiro, Martin Mottet, Jodie Muir, Mark A. Nelson, Alexander Ostrowski, Keith Parkinson, Anna Podedworna, Steve Prescott, Chris Rallis, Joshua Raphael, Richard Sardinha, Chris Seaman, Rudy Siswanto, Craig J Spearing, Ron Spencer, Leroy Steinmann, Zack Stella, Matt Stewart, Matt Stikker, David C. Sutherland, Autumn Rain Turkel, Paolo Vacala, Svetlin Velinov, Raoul Vitale, Jabari Weathers, Campbell White
\item \textbf{Cartographers:} Francesca Baerald, Jared Blando, CoupleOfKooks, Andrew Kolb, Dyson Logos, Damien Mammoliti, Mike Schley, Paolo Vacala
\item \textbf{Consultants:} Shawn Banerjee, Eytan Bernstein, Daniel Delgado, Sameer Joseph
\item \textbf{Producers:} Rob Hawkey (lead), Bill Benham, Siera Bruggeman, Vanessa Hoskins
\end{itemize}
\textit{This book is dedicated to Christopher Perkins for his decades of contributions to Dungeons \& Dragons.}
% Content: Unknown (dict)
% Content: Unknown (dict)

\end{document}
